\chapter{Introdução} % Main chapter title
\label{cap:introdução} % For referencing the chapter elsewhere, use Chapter~\ref{cap:introdução}

\textbf{Este guião de elaboração de relatório é apenas um guia não substituindo o diálogo necessário entre estudante e orientador para a melhor definição da estrutura e conteúdo do relatório em cada caso concreto.}

\textbf{O template foi adaptado de \cite{TMDEIThesisTemplate}, onde consta mais informação que pode servir de apoio à escrita deste documento. Devem consultar.}


A Introdução deve dar ao leitor a informação básica necessária por forma a facilitar o enquadramento dos objetivos, da abordagem seguida e dos resultados. 

A introdução começa com uma perspetiva geral do problema em estudo (na secção de enquadramento) em que, à medida que vai progredindo, deve ir fornecendo informação mais específica até se abordar a área em concreto tratada no relatório. 

Deve descrever, de forma sucinta, o problema em estudo e os objetivos. Deve também enunciar os principais métodos utilizados no trabalho, bem como a identificação clara dos contributos e aspetos inovadores da solução. A introdução deve terminar com a apresentação sucinta das secções que fazem parte do resto do relatório. 


\textbf{Nota: devem usar frases curtas; adotar o impessoal em vez do pessoal (e.g. adotou-se vs. adotei vs. adotamos); usar verbos conjugados no presente; evitar encadear verbos seguidos (e.g. “esta secção vai ser descrito” vs. “esta secção descreve”); usar voz passiva vs. ativa.}


%----------------------------------------------------------------------------------------

\section{Enquadramento/Contexto} 
\label{sec:enquadramento} %For referencing this section elsewhere, use Section~\ref{sec:enquadramento}

Enquadrar o problema no âmbito do projeto da organização proponente e descrever contexto e motivação do problema. 


%----------------------------------------------------------------------------------------

\section{Descrição do Problema}
O problema deve ser descrito de forma clara, de forma a que o leitor compreenda facilmente.


\subsection{Objetivos} 
Os objetivos devem ser descritos de forma detalhada devendo refletir a compreensão do estudante sobre o trabalho que foi desenvolvido.


\subsection{Abordagem} 
A aproximação preconizada para a solução do problema ou do tratamento do tema focado, onde estejam claras as contribuições previstas. O objetivo é identificar a abordagem adotada (usando referências bibliográficas) e como foi aplicada ao projeto. 


\subsection{Contributos} 
Devem ser apresentados os aspetos inovadores e de realce do trabalho, bem como a identificação dos benefícios trazidos para a organização (para a sociedade, para o ambiente,…). Devem ser identificadas as contribuições previstas (podendo usar estilo de apresentação por itens). 


\subsection{Planeamento do trabalho}  
Nesta subsecção deve ser apresentado o planeamento (cronograma, etc.) definido para a execução do projeto com indicação das tarefas, datas e principais fases (milestones).

\textbf{(Poderá ir para Anexo.)}


%----------------------------------------------------------------------------------------

\section{Estrutura do relatório}

Apresentação sucinta dos capítulos que fazem parte do relatório, descrevendo em poucos parágrafos o que cada um deles trata. 

Para além da Introdução, este documento contém mais X capítulos.

No Capítulo~\ref{cap:estado_da_arte}, é descrito o estado da arte e são apresentados trabalhos relacionados.

No Capítulo~\ref{cap:analise_desenho}, ...

\vspace{20mm} 

\textbf{(Este capítulo deverá ter no máximo 5 páginas.)}
